\documentclass[11pt]{article}

% Change "review" to "final" to generate the final (sometimes called camera-ready) version.
% Change to "preprint" to generate a non-anonymous version with page numbers.
\usepackage[review]{acl}

% Standard package includes
\usepackage{times}
\usepackage{latexsym}

% For proper rendering and hyphenation of words containing Latin characters (including in bib files)
\usepackage[T1]{fontenc}
% For Vietnamese characters
% \usepackage[T5]{fontenc}
% See https://www.latex-project.org/help/documentation/encguide.pdf for other character sets

% This assumes your files are encoded as UTF8
\usepackage[utf8]{inputenc}

% This is not strictly necessary, and may be commented out,
% but it will improve the layout of the manuscript,
% and will typically save some space.
\usepackage{microtype}

% This is also not strictly necessary, and may be commented out.
% However, it will improve the aesthetics of text in
% the typewriter font.
\usepackage{inconsolata}

%Including images in your LaTeX document requires adding
%additional package(s)
\usepackage{graphicx}
\usepackage{booktabs}
\usepackage{placeins}  % \FloatBarrier — keep tables/figures within their section

% If the title and author information does not fit in the area allocated, uncomment the following
%
%\setlength\titlebox{<dim>}
%
% and set <dim> to something 5cm or larger.

\title{How To Report Evaluations of your LLM: A Story of LLM/Human Misunderstandings}

% Author information can be set in various styles:
% For several authors from the same institution:
% \author{Author 1 \and ... \and Author n \\
%         Address line \\ ... \\ Address line}
% if the names do not fit well on one line use
%         Author 1 \\ {\bf Author 2} \\ ... \\ {\bf Author n} \\
% For authors from different institutions:
% \author{Author 1 \\ Address line \\  ... \\ Address line
%         \And  ... \And
%         Author n \\ Address line \\ ... \\ Address line}
% To start a separate ``row'' of authors use \AND, as in
% \author{Author 1 \\ Address line \\  ... \\ Address line
%         \AND
%         Author 2 \\ Address line \\ ... \\ Address line \And
%         Author 3 \\ Address line \\ ... \\ Address line}

\author{First Author \\
  Affiliation / Address line 1 \\
  Affiliation / Address line 2 \\
  Affiliation / Address line 3 \\
  \texttt{email@domain} \\\And
  Second Author \\
  Affiliation / Address line 1 \\
  Affiliation / Address line 2 \\
  Affiliation / Address line 3 \\
  \texttt{email@domain} \\}

%\author{
%  \textbf{First Author\textsuperscript{1}},
%  \textbf{Second Author\textsuperscript{1,2}},
%  \textbf{Third T. Author\textsuperscript{1}},
%  \textbf{Fourth Author\textsuperscript{1}},
%\\
%  \textbf{Fifth Author\textsuperscript{1,2}},
%  \textbf{Sixth Author\textsuperscript{1}},
%  \textbf{Seventh Author\textsuperscript{1}},
%  \textbf{Eighth Author \textsuperscript{1,2,3,4}},
%\\
%  \textbf{Ninth Author\textsuperscript{1}},
%  \textbf{Tenth Author\textsuperscript{1}},
%  \textbf{Eleventh E. Author\textsuperscript{1,2,3,4,5}},
%  \textbf{Twelfth Author\textsuperscript{1}},
%\\
%  \textbf{Thirteenth Author\textsuperscript{3}},
%  \textbf{Fourteenth F. Author\textsuperscript{2,4}},
%  \textbf{Fifteenth Author\textsuperscript{1}},
%  \textbf{Sixteenth Author\textsuperscript{1}},
%\\
%  \textbf{Seventeenth S. Author\textsuperscript{4,5}},
%  \textbf{Eighteenth Author\textsuperscript{3,4}},
%  \textbf{Nineteenth N. Author\textsuperscript{2,5}},
%  \textbf{Twentieth Author\textsuperscript{1}}
%\\
%\\
%  \textsuperscript{1}Affiliation 1,
%  \textsuperscript{2}Affiliation 2,
%  \textsuperscript{3}Affiliation 3,
%  \textsuperscript{4}Affiliation 4,
%  \textsuperscript{5}Affiliation 5
%\\
%  \small{
%    \textbf{Correspondence:} \href{mailto:email@domain}{email@domain}
%  }
%}

\begin{document}
\maketitle
\begin{abstract}
Abstract to be added.
\end{abstract}

\section{Introduction}
\label{sec:introduction}

% Imagine the same model evaluated on the same benchmark by two
% leaderboards. The two scores differ by more than thirty points. Neither
% leaderboard is wrong, and neither is hiding anything: each reports the
% score its pipeline produced. What is missing is the procedure that
% turned the model into the number --- which harness was used, which
% prompt template, which scoring mode, how many in-context examples.
% Without these, the two scores are not comparable, and the gap between
% them is not interpretable.
%
% Evaluation results travel through a chain of hands before a reader sees
% them. A model author runs an internal evaluation; a leaderboard
% re-evaluates the same checkpoint with its own harness; an aggregator
% scrapes the leaderboard and stores the score in its own schema; a
% downstream paper cites the aggregator. At each hop, information about
% how the score was produced is lost, dropped, or quietly reinterpreted.
% By the time the reader compares two numbers, the procedures behind
% them may share no more than a benchmark name.
%
% This paper audits one large slice of that chain. We ingest 29{,}331
% evaluation records from 11 public sources covering 5{,}672 models and
% 180 benchmarks into a single schema and ask, at each step, what
% information survives. We find that the configuration fields most
% responsible for score divergence --- harness identity, prompt
% template, temperature --- are essentially absent from public records,
% and that the structural overlap needed for independent cross-source
% verification is rare. We then look closely at a handful of cases where
% the same model received noticeably different scores and ask what would
% have been needed to explain the gap.
%
% \paragraph{Contributions.}
% \begin{itemize}
%   \item A cross-source audit of metadata coverage across 29{,}331
%   evaluation records from 11 public sources.
%   \item Three case studies of cross-source score divergence, each
%   diagnosed in terms of the missing information that would have made
%   the gap interpretable.
%   \item EvalSpec v0.1, a proposed machine-readable reporting
%   specification whose required fields are grounded in the audit.
%   \item Public release of the dataset, ingestion pipeline, schema, and
%   analysis code under CC-BY-4.0.
% \end{itemize}
%
% The rest of the paper is organised as follows. Section~\ref{sec:related}
% reviews prior audits of evaluation practice and introduces the EEE
% schema. Section~\ref{sec:materials} describes how the dataset was
% assembled and what it contains. Section~\ref{sec:quant} reports the
% quantitative analysis and Section~\ref{sec:cases} works through the
% case studies. Section~\ref{sec:discussion} draws out the implications
% and presents EvalSpec; Section~\ref{sec:conclusions} concludes.
%% END DRAFT ----------------------------------------------------------

Open LLM Leaderboard v2 (OLv2) reports Gemma-7B's score on Big-Bench
Hard (BBH) as 21.12. A collection of ArXiv papers, the Jamba, RecurrentGemma,
and MAmmoTH2 papers, among others reports BBH scores for the same model
between 55 and 60. The gap exceeds 35 percentage points. Neither side is
fabricating a result: each reports the score that its pipeline produced.
What is missing is the procedure. OLv2 runs BBH with lighteval using
log-likelihood scoring over answer choices, whereas the papers follow the original
BBH protocol of 3-shot chain-of-thought generation. Those are not the same measurement and a benchmark name is not sufficient
to make them so; Figure~\ref{fig:chain} traces this attrition for the
related Gemma-2-9B case (\S\ref{sec:sealion}).

This example is the rule, not the exception. Evaluation results travel
through a chain of hands before a reader encounters them: a model author runs
an internal sweep, a leaderboard re-evaluates the checkpoint, an aggregator
scrapes the leaderboard, and a downstream paper cites the aggregator. At
each stage, information about how the score was produced is lost, dropped,
or quietly reinterpreted. By the time two scores meet on a comparison table,
the procedures behind them may share nothing more than a benchmark name.
For that reason we assemble a dataset of 37{,}718 evaluation records from
12 sources covering 6{,}455 models and 1{,}289 benchmarks, and for a
10-source leaderboard subset of 28{,}755 records ask what configuration
information survives. The fields most responsible for score divergence are
missing from the leaderboard subset: prompt template and temperature have
zero valid coverage, and no leaderboard API exposes the evaluation harness at
all (where harness identity appears it is attached by our pipeline from
platform documentation, not returned by the source). A controlled factorial
experiment (\S\ref{sec:cases}) further shows that temperature itself drives
under 1\% of score variance. We then work through cases where the
same model received sharply different scores on the same benchmark and
diagnose the mechanism that produced each gap.

The case studies are the contribution, our extraction pipeline is the
instrument that found them. Each identifies a specific class of mechanism
through which benchmark scores become incomparable across sources without
either source committing an error: implementation choices left
underspecified by the benchmark, multiple legitimate measurements coexisting
under one name, and apparent ``independent'' evidence that shares a single
pipeline origin.

\textbf{Contributions.}
\begin{itemize}
    \item A cross-source audit of metadata coverage in the public evaluation ecosystem.
    \item Three case studies of cross-source score divergence, diagnosed end-to-end.
    \item EvalSpec v0.1, a proposed reporting specification grounded in the audit.
    \item Public release of the full dataset (37{,}718 records), ingestion
      pipeline, EvalSpec v0.1 vocabulary and reference validator, and analysis
      code under CC-BY-4.0, contributed back to the EvalEval shared task
      \citep{evaleval2026everyevalever} as a worked example of the EEE schema
      in cross-source use; the dataset and EvalSpec are designed to be
      extended by the community.
\end{itemize}

\section{Related Work}
\label{sec:related}
\subsection{Evaluation Mispractices and Prior Audits}

%% DRAFT --------------------------------------------------------------
% Several recent studies have shown that benchmark scores are sensitive
% to choices that public records do not usually expose. Alzahrani et
% al.~\cite{alzahrani2024} report score gaps of up to fifteen points on
% MMLU attributable to harness choice alone, holding the model and the
% benchmark fixed. Biderman et al.~\cite{biderman2024} extend the picture
% to prompt format, tokenisation, and generation configuration, and
% recommend that these be reported as part of any evaluation result. A
% complementary line of work examines the platforms that host these
% evaluations: Singh et al.~\cite{singh2025} audit Chatbot Arena and
% document score instability for the same checkpoint under different
% aliases, along with access asymmetries between providers. Data
% contamination~\cite{sainz2023} is a parallel reproducibility threat
% that documentation alone cannot resolve, but whose detection likewise
% benefits from richer per-run records.
%
% A separate body of work proposes documentation frameworks. Model
% Cards~\cite{mitchell2019} and Datasheets~\cite{gebru2021} established
% the norm of structured artefact descriptions for models and datasets.
% Three concurrent 2025 proposals --- BenchmarkCards~\cite{sokol2025},
% EvalCards~\cite{dhar2025}, and Eval Factsheets~\cite{bordes2025} ---
% extend this norm to evaluation, but each is narrative rather than
% machine-validated, and none is backed by an audit of how the public
% ecosystem currently records the fields they recommend. Table~\ref{tab:frameworks}
% summarises the differences.
%% END DRAFT ----------------------------------------------------------


Several recent studies have shown that benchmark scores are sensitive to
implementation choices that the public record rarely exposes.
\citet{alzahrani2024} report score gaps of up to fifteen percentage points on
MMLU attributable to harness choice alone, holding the model and benchmark
fixed. \citet{biderman2024} extend this idea to prompt format, tokenisation,
and generation configuration, and recommend that these choices be reported as
part of any evaluation result. A complementary line of work examines the
platforms that host evaluations: \citet{singh2025} audit Chatbot Arena and
document score instability for the same checkpoint under different aliases,
along with access asymmetries between providers. Data contamination
\citep{sainz2023} is a parallel reproducibility threat whose detection likewise depends on richer per-run records.

These studies establish that scores vary with protocol choices. This paper
traces the concrete mechanisms through which that variation enters the public
record as built: not which harness choice matters in the abstract, but which
specific decisions produce which specific gaps in our cross-source dataset. The missing fields are not exotic, they are quantities known at evaluation time that are simply not exposed downstream.

A separate body of work proposes structured documentation for evaluation
artefacts. Model Cards \citep{mitchell2019} and Datasheets for Datasets
\citep{gebru2021} established the norm of structured descriptions for models
and data. Three concurrent 2025 proposals extend this to evaluation pipelines:
BenchmarkCards \citep{sokol2025}, EvalCards \citep{dhar2025}, and Eval
Factsheets \citep{bordes2025}. Each proposes fields that overlap substantially
with those identified as absent in our audit. The key difference is that these
frameworks are prescriptive: they specify what should be recorded. Our work is
empirical first, we establish what the public ecosystem currently records,
identify which absences produce uninterpretable divergences, and use that
evidence to ground the field selection in EvalSpec~v0.1 (Section~\ref{sec:evalspec}).


\subsection{The EEE Schema}

%% DRAFT --------------------------------------------------------------
% Our analysis builds on the Every Eval Ever (EEE) schema, a unified
% JSON format for storing evaluation records from heterogeneous sources.
% Each record captures one model's result on one benchmark from one
% source, organised into four blocks: \texttt{model\_info} describes
% the model under evaluation; \texttt{eval\_library} describes the
% harness used to produce the score; \texttt{source\_metadata} describes
% the source of the record and its provenance; and
% \texttt{evaluation\_results} contains the score itself along with
% per-run generation configuration. Table~\ref{tab:eee-fields} lists the
% fields most relevant to this paper.
%
% A unified schema is the right abstraction for cross-source comparison
% because the alternative --- comparing scores across sources whose
% formats and conventions differ --- requires implicit alignment that
% nobody verifies. The schema makes the alignment explicit and turns
% missing information from an invisible gap into a recordable absence:
% a field that no source populates is visibly empty in every record,
% rather than silently inferred away.
%% END DRAFT ----------------------------------------------------------

Our analysis builds on the Every Eval Ever (EEE) schema, introduced by the
EvalEval Coalition \citep{evaleval2026everyevalever}, which articulates the
cross-framework score-comparability problem this paper investigates (citing
the HELM--EleutherAI gap on LLaMA~65B MMLU, 0.637 vs.\ 0.488); our case
studies trace the specific mechanisms behind such gaps in the public record.
EEE is a unified JSON format for storing evaluation records from heterogeneous
sources. Each record captures one model's result on one benchmark from one
source, organised around four main objects. \texttt{model\_info} records the
model name and unique identifier, with an optional \texttt{developer} field.
\texttt{eval\_library} records the harness name and version.
\texttt{source\_metadata} records the source organisation, source type
(\texttt{documentation} or \texttt{evaluation\_run}), and the evaluator's
relationship to the model (\texttt{first\_party}, \texttt{third\_party},
\texttt{collaborative}, or \texttt{other}).
\texttt{evaluation\_results} is an array of result objects, each containing a
\texttt{metric\_config} block (metric name and scoring details) and a
\texttt{score\_details} block (score value and uncertainty); a
\texttt{generation\_config} block (generation arguments such as temperature and
prompt template) is optional.
A unified schema is the right abstraction for cross-source comparison because
the alternative requires implicit alignment that is never written down and never verified. The schema makes alignment explicit and transforms missing information from an invisible gap into a recordable absence: a field that no source populates appears as empty in every record, rather than being silently inferred.
It is precisely the fields that are visibly empty across all records that
motivate this paper. The dataset and code are released under CC-BY-4.0.

\section{Materials}
\label{sec:materials}

\subsection{How We Gathered the Data}
\label{sec:pipeline}

%% DRAFT --------------------------------------------------------------
% We ingest evaluation records from eleven public sources: eight
% leaderboards (Open LLM Leaderboard v2, AlpacaEval 2, Chatbot Arena,
% BigCodeBench, EvalPlus, BFCL, SWE-Bench, WildBench, MT-Bench), the
% HuggingFace model-card YAML model-index, and Papers With Code. Each
% source is handled by a dedicated converter that maps its native
% format onto the EEE schema; records that fail JSON-schema validation
% are dropped, and duplicates are deduplicated on the composite key
% \texttt{(model\_id, benchmark, source)}.
%
% The pipeline is not uniform in how much it trusts each source. We
% distinguish three levels.
%
% \paragraph{Deterministic.} Schema validation and ID normalisation are
% rule-based and do not depend on any heuristic. Model identifiers are
% lowercased; \texttt{-hf} suffixes and shot annotations are stripped;
% benchmark names are normalised to a canonical form. These operations
% are reversible and produce no false positives.
%
% \paragraph{Regex-based and fragile.} Papers With Code does not expose
% a structured API; results are scraped from HTML tables whose layout
% varies by paper. Score values, model identifiers, and benchmark names
% are extracted by regex. We treat this layer as the least reliable
% part of the pipeline and audit it explicitly in
% Section~\ref{sec:pwc-audit}.
%
% \paragraph{External knowledge.} Several sources do not expose the
% harness they use in their data, but the harness is known from
% platform documentation (for example, OLv2 uses lighteval; AlpacaEval
% uses alpaca\_eval). We record these inferred values in a separate
% \texttt{inferred\_harness} field rather than in
% \texttt{eval\_library.name}, so that self-reported and inferred
% values remain distinguishable downstream.
%
% Cross-source collisions are detected after normalisation: records are
% grouped by \texttt{(model\_id, benchmark)}, and every pair of records
% from distinct sources within a group is emitted as a collision pair
% with its score delta and metadata-match annotations. The full
% procedure is described in Appendix~\ref{app:collision}.
%
% We deliberately abstain from two operations that would otherwise be
% tempting. We do not use a language model to re-extract scores from
% paper PDFs or HTML, because doing so would introduce a second layer of
% errors that we could not separate from the errors already present in
% the sources. We do not manually correct individual records, because
% the audit's value depends on reporting what the public ecosystem
% actually contains, not what it would contain after curation.
%% END DRAFT ----------------------------------------------------------
The dataset is assembled from three structurally different classes of
source. Leaderboards are taken as authoritative: each leaderboard fetcher
downloads its data at runtime and emits a schema-conformant record per
model$\times$benchmark cell, with the n-shot count and chain-of-thought
flag attached from the leaderboard's fixed protocol. OLv2 dominates and
exposes only model identifiers, scores, and a submission date; its
\texttt{eval\_library.name} is left as \texttt{unknown}, and no field for
evaluation harness, prompt template, scoring mode, or temperature is
returned by the API. HuggingFace Model Cards are parsed for the YAML
\texttt{model-index} block, capturing first-party results that model authors
chose to document. The third source is ArXiv paper extraction: for 100
papers, the HTML rendering is fetched and an LLM (Claude Haiku via
OpenRouter) reads paper prose into the EEE schema, returning per-benchmark
metadata that may include n-shot count, temperature, prompt template,
chain-of-thought flag, and harness name when explicitly stated. The LLM may
hallucinate or misattribute values, so we treat these metadata fields as
unverified at scale.


\subsection{What We Gathered}
\label{sec:dataset}

%% DRAFT --------------------------------------------------------------
% The resulting dataset contains 29{,}331 evaluation records covering
% 5{,}672 unique models and 180 unique benchmarks across the eleven
% sources. Table~\ref{tab:sources} reports per-source counts.
%
% The dataset is shaped by three sources whose properties are worth
% naming upfront, because they recur throughout the analysis. Open LLM
% Leaderboard v2 contributes 26{,}790 records, or 91\% of the total;
% any aggregate figure we report is dominated by it unless explicitly
% ablated. HuggingFace Model Cards are the opposite shape: 207 records,
% only 11 unique models, but 156 distinct benchmarks --- a few model
% authors who report many benchmarks each. Papers With Code is the
% messiest source: of 1{,}350 raw records, 774 were removed as
% extraction artefacts or duplicates, leaving 576 in the clean dataset.
% The remaining sources contribute benchmark diversity but limited
% model overlap.
%
% Figure~\ref{fig:coverage-matrix} visualises the source~$\times$~benchmark
% coverage matrix. Most cells are empty: each source evaluates a
% relatively small subset of the benchmark universe, and the subsets
% overlap less than one might expect. This sparsity is the precondition
% for the structural fragmentation we report in Section~\ref{sec:quant}.
%% END DRAFT ----------------------------------------------------------


The full dataset contains 37{,}718 evaluation records covering
6{,}455 unique models and 1{,}289 unique benchmarks across 12 sources.
Table~\ref{tab:sources} reports per-source counts.

The 10 leaderboard sources contribute 28{,}755 records across 5{,}301 models
and 174 benchmarks. Open LLM Leaderboard v2 (OLv2) dominates, contributing
26{,}790 records (93.2\%) across 4{,}465 models on 6 benchmarks.
HuggingFace Model Cards present the opposite shape: 207 records from only
11 unique models but spanning 156 distinct benchmarks, a handful of model
authors documenting many benchmarks each. Chatbot Arena, SWE-Bench, and
MT-Bench each report one record per model, reflecting their pairwise or
task-specific evaluation structure. The remaining leaderboard sources
(AlpacaEval 2, BFCL, BigCodeBench, EvalPlus, WildBench) contribute
benchmark-specific coverage within their respective domains.

The ArXiv extraction source contributes a further 8{,}387 records from
100 papers (counting only the canonical per-benchmark LLM extraction; audit
re-extractions are not double-counted). Papers With Code contributes a further
576 records but is excluded from all statistical analyses due to systematic
scraping artefacts (parameter counts parsed as scores). The ArXiv source is
excluded from the headline leaderboard-subset
metadata coverage analysis (Figure~\ref{fig:coverage}) because its metadata
fields were LLM-populated rather than read from a structured API; coverage
figures would otherwise be inflated by unverified values. The 10-source
leaderboard subset measures what platforms expose in machine-readable form,
which is the relevant measure of documentation practice. The ArXiv source
instead serves the case studies in Section~\ref{sec:cases}, where it
provides the cross-source comparison point: a score reported by a
leaderboard is compared against the same model$\times$benchmark cell as
reported in the original paper.

\begin{table}[t]
\centering
\small
\begin{tabular}{lrrr}
\toprule
Source & Records & Models & Benchmarks \\
\midrule
Open LLM Leaderboard v2 & 26{,}790 & 4{,}465 & 6 \\
AlpacaEval 2 & 446 & 223 & 2 \\
BFCL & 327 & 109 & 3 \\
BigCodeBench & 280 & 156 & 2 \\
Chatbot Arena & 218 & 218 & 1 \\
EvalPlus & 214 & 125 & 2 \\
HF Model Card & 207 & 11 & 156 \\
WildBench & 122 & 61 & 2 \\
SWE-Bench & 117 & 117 & 1 \\
MT-Bench & 34 & 34 & 1 \\
\midrule
\textit{Leaderboard subtotal} & \textit{28{,}755} & \textit{5{,}301} & \textit{174} \\
\midrule
ArXiv papers (LLM extraction) & 8{,}387 & 783 & 1{,}134 \\
Papers With Code & 576 & 383 & 16 \\
\midrule
\textbf{Total (12 sources)} & \textbf{37{,}718} & \textbf{6{,}455} & \textbf{1{,}289} \\
\bottomrule
\end{tabular}
\caption{Per-source record, model, and benchmark counts. The leaderboard
subtotal defines the 10-source subset used for metadata coverage analysis;
the ArXiv extraction source provides the cross-source comparison point for
the case studies. Model and benchmark totals do not sum because many appear
in multiple sources.}
\label{tab:sources}
\end{table}

\section{Quantitative Analysis}
\label{sec:quant}

\subsection{Human Annotation of Automatic Score Extractions}
\label{sec:score-ann}

\paragraph{Score extraction accuracy (ground-truth check).}
We manually verified all score extractions across 10 held-out ArXiv papers
(1{,}619 entries spanning DataComp-LM, DeepSeek LLM, Instella, Kimi-Audio,
Llama-Embed-Nemotron-8B, LLM360, Mistral~7B, OpenELM,
Phi-4-reasoning-vision, and Sailor), confirming 100\% agreement between
extracted and published scores. Of these, 34 entries (21 from
Phi-4-reasoning-vision and 13 from Mistral~7B) report round numbers
(e.g.\ 50, 84, 1{,}012) that appear in the paper as integers rather than the
floats stored in the JSON; these are considered correct since the numeric
values are identical. Per-paper row counts are reported in
Table~\ref{tab:score-ann-app} (Appendix~\ref{app:audit}). This is a ground-truth accuracy check on the numeric extraction
step: it establishes that when the pipeline places a value in the score
field, that value matches the paper. It does not verify that the extraction
correctly captures metadata (harness, n-shot, metric variant), which is a
harder problem.

\paragraph{Structural presence check.}
Across the 8{,}387 LLM-extracted ArXiv records,
8{,}192 (97.7\%) carry a non-null score value. This is a JSON-level
structural check, which confirms that the extraction pipeline placed a value
in the score field and is not a verification that the value is semantically
correct. 

\subsection{Human Annotation of Automatic Harness Identification}
\label{sec:harness-ann}
We assembled two balanced samples from the ArXiv extraction source: five
papers that name an evaluation harness (\emph{with-harness}: BLOOM, Llemma,
Sheared LLaMA, Aya 23, Nemotron-4 15B) and five that do not (\emph{without-harness}:
Code Llama, Mamba, Reka, Jamba, and a Llama-2-related survey paper),
covering 1{,}633 extracted (model, benchmark) rows in total. For each paper, the
ArXiv HTML was retrieved and the pipeline's per-row harness assignment was
verified against the source text under a strict rubric: the pipeline's
\texttt{eval\_library.name} value for a row is \textsc{true} only if it
correctly names the tool the paper says was used to produce that specific
score (with canonical equivalents like ``eval-harness'' and ``Language
Model Evaluation Harness'' accepted as the same library, but a paper's
distinctly-named fork required as a verbatim match). Methodology is in
Appendix~\ref{app:harness}.

On the with-harness sample, per-row precision was 704/817 = 86.2\%. Three
of five papers were at 100\% precision (BLOOM 436/436 with the
\emph{Prompted Language Model Evaluation Harness} fork name correctly
extracted on all 91 SuperGLUE rows; Sheared LLaMA 138/138; Aya 23 58/58).
Two papers exposed a real pipeline failure mode: when a paper attributes
a named harness to a broad group via a phrase like ``across all
aforementioned tasks'' (Nemotron-4 15B §3.1) or ``for open access
models'' (Llemma §3.1), the LLM under-extracts and applies the
attribution to only the benchmarks explicitly listed in the same
sentence, missing the broader group. Nemotron-4 15B precision was 6/96
(only the 6 commonsense-reasoning rows tagged, with the remaining 90
math/code/multilingual rows in the same \emph{aforementioned} group left
as \texttt{unknown}); Llemma was 66/89 (the LLM caught some open-access
model rows but not all).

On the without-harness sample, abstention quality was 816/816 = 100\%.
The pipeline returned \texttt{unknown} for every row across all five
papers, and manual inspection confirmed that none of the five names a
public harness for its own evaluations. Three of the five (Mistral 7B's
``our own evaluation pipeline'', Phi-3's ``Microsoft internal tool'',
Reka's Reka-Core-as-judge internal stack) are \emph{named-but-not-public}
cases, the paper acknowledges a proprietary tool without giving it a
public identifier; the remaining two (Code Llama, Mamba) name nothing at
all.


\subsection{Results and Discrepancies}
\label{sec:discrepancies}

\paragraph{Metadata coverage.}
Figure~\ref{fig:coverage} reports coverage for five configuration fields
across both the full 37{,}718-record dataset and the 10-source
28{,}755-record leaderboard subset. Coverage counts a field as populated
whenever the EEE record carries a non-``unknown'' value, irrespective of
whether that value was returned by the source API, attached by the
pipeline from platform documentation, or extracted from paper prose by
the LLM extractor.

\begin{figure}[t]
\centering
\includegraphics[width=\linewidth]{./figures/coverage_bars.pdf}
\caption{Metadata coverage across five configuration fields, full
  dataset (37{,}718 records) vs.\ the 10-source leaderboard subset
  (28{,}755 records). Prompt template and temperature have zero valid
  coverage in the leaderboard subset. Harness identity shows 93\% in the
  subset, but this is entirely pipeline-attached: OLv2's API does not
  expose a harness field, and the pipeline writes \texttt{lm\_eval} to
  every OLv2 record from platform documentation. Full numerical values
  are in Table~\ref{tab:coverage-app}
  (Appendix~\ref{app:coverage}).}
\label{fig:coverage}
\end{figure}


The headline pattern is that the fields most responsible for score
divergence are either absent from the leaderboard subset
(\texttt{prompt\_template} and \texttt{temperature} at 0\%) or, in the
case of harness identity, populated only because the pipeline attaches
\texttt{lm\_eval} to every Open-LLM-Leaderboard~v2 record from platform
documentation (\S\ref{sec:pipeline}). No leaderboard source we examined
returns the harness identity via its API; the 93\% leaderboard-subset
harness coverage in Figure~\ref{fig:coverage} is the pipeline's auto-attachment, not source-recorded
data. As Section~\ref{sec:bbh} documents, OLv2 actually runs
\emph{lighteval}, and the discrepancy between the pipeline-attached
\texttt{lm\_eval} label and the actual \emph{lighteval} runtime is itself
an instance of the same problem: the API exposes nothing, so any
attribution downstream is inferred.

\section{Case Studies}
\label{sec:cases}

The case studies are organised into three groups corresponding to
three distinct mechanisms through which benchmark scores become incomparable
across sources.

To quantify the relative importance of methodological factors before
examining specific failure modes, we ran a controlled factorial experiment
crossing four instruction-tuned models (Qwen2.5-14B, Qwen2.5-7B,
Llama-3.1-8B, Mistral-7B-v0.3) against three benchmarks (BBH, GSM8K, MMLU)
under three prompt formats (plain few-shot, instruction-formatted,
chain-of-thought), two temperature settings (0.0 and 0.7), and two n-shot
counts per benchmark (0-shot plus 3-shot for BBH, 0-shot plus 5-shot for
GSM8K and MMLU), with three random seeds per cell; MMLU was restricted to
temperature 0.0 and a single seed. The resulting dataset comprises 306
evaluations. Prompt format dominates on GSM8K: at 5-shot, it accounts
for approximately 27\% of variance (\(F = 13.0\), \(p < 0.001\); 24
cells per format). MMLU's reduced 9-cell design at 5-shot is
descriptively consistent (\(\eta^2 \approx 24.5\%\)) but underpowered
(\(F = 0.98\), \(p = 0.43\)); the per-format means agree with the GSM8K
pattern but cannot independently confirm it. Temperature contributes
under 1\% on GSM8K (\(\eta^2 = 0.25\%\)); MMLU was run at a single
temperature, so the temperature comparison is not available. The n-shot
count's contribution is benchmark-dependent: negligible on MMLU
(\(\eta^2 \approx 1\%\)), but on GSM8K inseparable from the answer parser,
because 0-shot runs collapse to zero when the models omit the delimiter the
extractor expects (\S\ref{app:exp-grid}). Full
per-cell results and per-benchmark ANOVAs are in
Appendix~\ref{app:experiments}. The three case studies below trace this
pattern to specific, mechanistically distinct interactions between format
choices and evaluation infrastructure.

\begin{figure*}[t]
\centering
\includegraphics[width=\textwidth]{./figures/case_studies_unified_belebele_S3.pdf}
\caption{Three mechanisms of score incomparability
(\S\S\ref{sec:protocol}--\ref{sec:convergence}).
Each column compares two protocol variants under the same benchmark label
(blue panel = top configuration, orange = bottom;
$\Delta =$ bottom $-$ top in percentage points).
\textbf{Left (\S\ref{sec:gsm8k}):} GSM8K on Qwen2.5-14B-Instruct, 5-shot,
lm-evaluation-harness 0.4.11~--- plain vs.\ chain-of-thought prompt format.
CoT elicits prose without the \texttt{\#\#\#\#} delimiter;
parser records 0. $\Delta = -56.0$ pp.
\textbf{Centre (\S\ref{sec:belebele}):} Belebele on Mistral-7B.
Reka paper \citep{team2024reka}: 0-shot, multilingual average (150 languages),
32.80. Falcon2 report \citep{malartic2024falcon2}: 5-shot, English only, 79.42.
Falcon2's internal 0-shot English score (32.48) isolates the n-shot effect.
$\Delta = +46.6$ pp.
\textbf{Right (\S\ref{sec:sealion}):} BBH on Gemma-2-9B across five sources.
OLv2 and SEA-LION share a pipeline origin (exact 2-decimal match, 34.10);
OLMoE/SmolLM2/RecurrentGemma run 0-shot CoT independently (68.20--69.00).
$\Delta = +34.6$ pp (cluster gap).}
\label{fig:casestudies-unified}
\end{figure*}

\subsection{Case Study 1: Errors from Undocumented Scoring Mode}
\label{sec:protocol}

The mechanism in this case study: the harness that runs a benchmark makes
implementation choices that the benchmark specification leaves underspecified,
and different choices produce measurements that cannot be compared on the same
scale, even when both are correctly executed.

\begin{figure}[t]
\centering
\includegraphics[width=\linewidth]{./figures/case1_parser_flow_S3.pdf}
\caption{Parser-flow diagram for Case Study 1 (\S\ref{sec:gsm8k}).
The same Qwen2.5-14B-Instruct model output is routed through the same
\texttt{\#\#\#\#} delimiter regex under two prompt formats.
Plain few-shot output contains the delimiter; the regex matches and
\texttt{exact\_match}$=1$. Chain-of-thought output contains step-by-step
prose without the delimiter; \texttt{filtered\_resps}$=$\texttt{`[invalid]'}
and \texttt{exact\_match}$=0$.
Aggregate: 58.5\% (plain) vs.\ 2.5\% (CoT), $\Delta = -56.0$ pp.}
\label{fig:case1-parser}
\end{figure}


\subsubsection{BBH: Log-Likelihood Scoring vs.\ Chain-of-Thought Generation}
\label{sec:bbh}

The Big-Bench Hard (BBH) benchmark was introduced by Suzgun et
al.~\cite{suzgun2023} with a specific evaluation protocol: 3-shot chain-of-thought
generation, where the model produces free-form reasoning text and the final
answer is extracted from that text. This protocol is the standard that ArXiv
papers reporting BBH scores follow, and it is the protocol implied by
the benchmark's leaderboard in the literature.

OLv2 evaluates BBH differently. It uses \emph{lighteval} with
log-likelihood scoring: for each question, the model is presented with
each candidate answer and the answer receiving the highest token log-probability
is selected. No text generation occurs, and no chain-of-thought reasoning
is elicited. OLv2's API does not expose the evaluation harness; the
harness field in our dataset is empty for all OLv2 records. The actual
harness (\emph{lighteval}) and its scoring mode (log-likelihood over answer
choices) are confirmed by platform documentation and harness source-code
inspection.

The consequences are visible in the data. Per-source BBH scores for Gemma-7B are in Table~\ref{tab:bbh-app} (Appendix~\ref{app:casetables}), which shows
OLv2's Gemma-7B BBH score of 21.12 against a cluster of ArXiv papers
reporting 55.10 (Jamba paper, RecurrentGemma paper), 57.40 (MAmmoTH2 paper),
and 59.60 (reported as ``BigBench-Hard'', which does not currently match the
``BBH'' benchmark alias in our normaliser). The gap exceeds 35 percentage
points. The same pattern holds for Gemma-2-2B (OLv2: 11.76; Instella paper:
40.80) and for MAmmoTH2-7B-Plus (OLv2: 18.93; MAmmoTH2 paper: 63.10).

A controlled scoring-mode comparison anchors this diagnosis. E2
(Appendix~\ref{app:exp-e2}) rules out a generic scoring-mode effect: the
log-likelihood vs.\ generation gap on MMLU is at most 0.28 pp across three
paired instruction-tuned models, two orders of magnitude smaller than the
BBH gap. The same infrastructure drives MMLU-Pro and any other benchmark
OLv2 evaluates via lighteval log-likelihood; cross-source comparisons
involving those scores should be treated with the same caution.

\subsubsection{MATH Lvl 5: Answer Parser Replacement}
\label{sec:math}

OLv2 evaluates the MATH Level 5 subset (1{,}324 problems, 5-shot) using
a SymPy-based answer parser. In February 2025, HuggingFace replaced this
parser with Math-Verify and re-evaluated all submitted
models~\cite{kydlicek2025mathverify}.

The old parser failed on \textbackslash{}boxed\{\} outputs not matching a
specific format string, mishandled sets, matrices, and interval notation via
SymPy, and did not implement numerical-equivalence checks (e.g., treating
$1/3$ and $0.333\ldots$ as different answers). The transition had measurable
effects: an average gain of $+4.66$ percentage points across all 3{,}751
re-evaluated models; $+8.27$ points on the Algebra subset; $+6.93$ points
on the Prealgebra subset; and extreme individual cases up to approximately
90 points on those subsets.

Any cross-source comparison that places a pre-transition OLv2 MATH Lvl 5
score alongside a post-transition score, or alongside a score from an
external paper that uses a correct parser, is comparing incomparable
quantities. OLv2 exposes a model submission date but not an
evaluation-execution timestamp, so it is impossible to distinguish pre- and
post-transition scores from the record alone: a model submitted in 2024
now carries a post-transition score despite predating the parser change by
months.

\subsubsection{GSM8K: Chain-of-Thought Format and Answer-Parser Failure}
\label{sec:gsm8k}

The answer-parser sensitivity documented in Section~\ref{sec:math} extends
beyond OLv2's SymPy transition to any harness that applies a format-specific
extractor to a model output shaped by the prompt. Within the factorial
grid, the Qwen2.5-14B-Instruct $\times$ GSM8K $\times$ 5-shot slice (E1,
Appendix~\ref{app:experiments}) illustrates this directly: the model is
evaluated under three prompt formats using lm-evaluation-harness 0.4.11,
with three seeds at temperature 0.0 (Table~\ref{tab:e1-body}).


\begin{table}[ht]
\centering
\small
\begin{tabular}{lr}
\toprule
Prompt format & Score \\
\midrule
\texttt{plain} (direct-answer few-shot)      & 58.5\% \\
\texttt{instruct} (system-prompt-formatted)  & 57.0\% \\
\texttt{cot} (chain-of-thought few-shot)     & ~~2.5\% \\
\bottomrule
\end{tabular}
\caption{E1: Qwen2.5-14B-Instruct on GSM8K, 5-shot, temperature 0.0,
\texttt{limit=200}. Identical across seeds 42/123/7 (deterministic
decoding). Full E1 design in Appendix~\ref{app:exp-e1}.}
\label{tab:e1-body}
\end{table}

The chain-of-thought format reduces measured accuracy by $\approx$56 percentage
points. The mechanism is parser failure: the CoT prompt elicits extended
reasoning text, but the harness extracts the final answer by searching for the
standard GSM8K delimiter (\texttt{\#\#\#\#~\{answer\}}). When the model
produces multi-step reasoning followed by a summary, the regex fails to locate
the delimiter and records zero for the item. The model's reasoning is not
incorrect; the parser cannot read the output format the prompt instructed it
to produce.

This is the same structural failure as the Math-Verify case: the score
reflects the parser's assumptions about output format, not the model's
mathematical capability. A reader comparing a CoT-prompted score with a
plain-prompted score on GSM8K, or a score from one harness whose extractor
handles CoT output against one that does not, is comparing incommensurable
quantities under an identical benchmark label.

\subsection{Case Study 2: Differences from Implementation Inconsistency}
\label{sec:onenametwo}

The mechanism in this case study: unlike Case Study 1 (Section~\ref{sec:protocol}),
where a single benchmark specification is implemented differently by
different harnesses, here multiple distinct specifications legitimately
coexist under the same benchmark name. A reader comparing two scores labeled
with the same name may be comparing measurements of genuinely different things.

\subsubsection{TruthfulQA: Generation-Based Evaluation at Different N-Shot}
\label{sec:truthfulqa}

TruthfulQA~\cite{lin2022truthfulqa} includes both multiple-choice metrics
(MC1, MC2) and a generation-based metric: the percentage of answers that are
both truthful and informative, judged by a GPT-based classifier following the
original paper's protocol. These metrics are not interchangeable, but they operate
on different model outputs and different measurement scales.

Our dataset contains records for Llama-2-Chat-7B at 57.04 (one ArXiv paper)
and 26.30 (multiple ArXiv papers), all labeled ``TruthfulQA''. Verification
of the source papers reveals that \emph{both} scores use the generation-based
metric (\% truthful and informative), not MC1 or MC2:

\begin{itemize}
  \item The 57.04 figure comes from the Code Llama paper
    \citep{roziere2023codellama} (arXiv:2308.12950), which evaluates
    Llama-2-Chat-7B as a baseline. The paper uses a \textbf{6-shot} prompt
    containing six random QA pairs in InstructGPT format, with a GPT-3-based
    judge. This methodology follows Lin et al.\ (2022) and is identical in
    metric to the Llama 2 paper's evaluation.
  \item The 26.30 figures come from papers using the Tülu evaluation suite
    \citep{ivison2023camels}, as exemplified by the OLMo paper
    \citep{groeneveld2024olmo} (arXiv:2402.00838). The Tülu suite applies the
    same generation-based metric at \textbf{0-shot}, but with a different judge:
    a Llama-2-7B fine-tuned classifier rather than the GPT-3-based judge used
    in the Code Llama evaluation.
\end{itemize}

The gap (57.04 vs.\ 26.30) therefore reflects two confounded differences: shot
count (6-shot vs.\ 0-shot) and judge implementation (GPT-3 vs.\ a Llama-2-7B
classifier). Both factors inflate the Code Llama score relative to the Tülu
score. Neither source paper records the shot count or the judge model in the
metric label; both simply report ``TruthfulQA''. Without explicit
\texttt{n\_shot} and judge fields in the record, the 30-point gap is
uninterpretable from the label alone. We isolate the judge factor directly in a
controlled experiment (Appendix~\ref{app:exp-judge}): holding a fixed set of
model responses constant and varying only the judge model shifts mean scores by
up to 1.35 points on a 10-point scale and can reverse the ranking of two
systems, confirming that judge identity alone---independent of shot count---is
enough to make two ``TruthfulQA'' scores incomparable.

\subsubsection{Belebele: N-Shot and Language Scope Confounded}
\label{sec:belebele}

Belebele~\cite{bandarkar2023belebele} is a multilingual reading-comprehension
benchmark. Two evaluation dimensions that should be separately labelled are
often not: the language scope (English-only vs.\ multilingual average) and the
n-shot count (which has a substantial effect on accuracy for this benchmark).

Our dataset records Mistral-7B receiving 32.80 from one ArXiv paper and
79.42 from another, both labeled ``Belebele''. Verification of the source
papers reveals that these two scores differ on both dimensions:

\begin{itemize}
  \item The 32.80 figure comes from the Reka paper \citep{team2024reka}
    (arXiv:2404.12387), which evaluates Belebele at \textbf{0-shot} as a
    \textbf{multilingual average} across 150 language variants. The evaluation
    was run by the Reka authors themselves (marked with a dagger in the paper).
  \item The 79.42 figure comes from the Falcon2 technical report
    \citep{malartic2024falcon2} (arXiv:2407.14885), which evaluates Belebele
    at \textbf{5-shot} on \textbf{English only}.
\end{itemize}

Crucially, the Falcon2 paper also reports a 0-shot English score for
Mistral-7B of 32.48,  nearly identical to Reka's 0-shot multilingual average
of 32.80. This internal comparison within the Falcon2 paper isolates the
n-shot effect: going from 0-shot to 5-shot on English Belebele yields the
full 47-point gain (32.48 $\to$ 79.42), while language scope alone (English
vs.\ multilingual average at 0-shot) contributes almost nothing. The dominant
driver of the 46.6-point gap between the two labeled scores is therefore the
n-shot count, not the language scope.

Neither record carries an \texttt{n\_shot} label or a language-scope field.
``Belebele 32.80'' and ``Belebele 79.42'' are not comparable measurements of
the same thing; they are two different evaluations of Mistral-7B that happen
to share a benchmark name.

\subsubsection{HellaSwag: Scoring Mode and N-Shot Confounded}
\label{sec:hellaswag}

HellaSwag is a sentence-completion benchmark for commonsense NLI. The benchmark
label conceals two evaluation dimensions at once: scoring mode and n-shot count.

Our dataset records Gemma-7B receiving 49.80 from one ArXiv paper and 81.20
from another, both labeled ``HellaSwag''. Verification of the source papers
reveals that the scores differ on both dimensions:

\begin{itemize}
  \item The 49.80 figure comes from the Phi-3 technical report
    \citep{abdin2024phi3} (arXiv:2404.14219), which evaluates Gemma-7B as a
    comparison baseline at \textbf{5-shot} using a \textbf{generation-based}
    pipeline: the model generates free text and the answer is extracted from
    its output (temperature~0, direct-answer format).
  \item The 81.20 figure comes from the Gemma technical report
    \citep{gemmateam2024gemma} (arXiv:2403.08295, Table~6), which reports
    Gemma~7B at \textbf{0-shot} using \textbf{log-likelihood} scoring over
    the four candidate continuations.
\end{itemize}

The $\approx$31-point gap is too large to be explained by the shot-count
difference alone: on a log-likelihood benchmark, varying n-shot typically shifts
scores by a few points at most, a bound we confirm directly in our factorial
grid, where MMLU's 0-shot and 5-shot log-likelihood means differ by 2.6 pp
($\eta^2 \approx 1\%$; Appendix~\ref{app:exp-grid}). The primary driver is the scoring mode.
Generation-based evaluation of HellaSwag, asking the model to produce the
correct continuation as free text, is structurally harder than log-likelihood
evaluation, where the model ranks pre-specified continuations by token
probability. Neither paper labels the scoring mode or the n-shot count in the
metric label; both simply report ``HellaSwag''. Without explicit
\texttt{scoring\_mode} and \texttt{n\_shot} fields in the record, the gap is
indistinguishable from a genuine capability difference.

\subsection{Case Study 3: Differences from Using a Different Harness}
\label{sec:convergence}

The mechanism in this case study: two sources reporting the same score to
multiple decimal places do not necessarily reflect two independent
measurements. They can reflect a single measurement that was rereported,
making them functionally a single data point disguising as two.

\subsubsection{SEA-LION and Gemma-2-9B on BBH}
\label{sec:sealion}

OLv2 reports Gemma-2-9B's BBH score as 34.10. The SEA-LION paper
(\citealt{ng2025sealion}; arXiv:2504.05747, IJCNLP-AACL 2025) reports the
same model's BBH score as 34.10, an exact match to two decimal places.
The SEA-LION paper lists its benchmark sources as ``SEA-HELM Leaderboard,
SEACrowd, Open LLM Leaderboard''. An exact two-decimal match of this kind is
consistent with either of two mechanisms: direct importation of the OLv2
score for a competitor model, or independent re-evaluation using identical
infrastructure (the same lighteval configuration and task suite that OLv2
runs). The EEE record cannot distinguish between them. Our harness-version
control (Appendix~\ref{app:exp-e3}) shows that independent reruns of the same
harness vary by up to 1.3 pp across seeds, so an \emph{exact} two-decimal match
is far more consistent with a re-imported score than with an independent
re-evaluation, which would almost never agree to the second decimal. In either
case, the two records share a pipeline origin and are not independent
measurements of Gemma-2-9B's BBH capability. Figure~\ref{fig:chain} visualises the
attrition that produces this collision: the fields that would distinguish
replication from re-import \texttt{harness},
\texttt{prompt\_template\_hash}, \texttt{evaluation\_timestamp} were
dropped before the records reached the comparison table.

Three other ArXiv papers (OLMoE, SmolLM2, RecurrentGemma) report
Gemma-2-9B's BBH score as 68.20, 69.00, and 69.00 respectively. These
papers used 0-shot chain-of-thought generation, a different protocol from
OLv2's 3-shot lighteval log-likelihood, explaining the $\approx$35-point
gap. The same harness-mismatch mechanism from Section~\ref{sec:bbh}
applies here as a compounding factor. All five sources and their scores
are listed in Table~\ref{tab:sealion-app} (Appendix~\ref{app:casetables}).


\section{Discussion}
\begin{figure*}[t]
\centering
\includegraphics[width=\textwidth]{./figures/eval_record_attrition_chain.pdf}
\caption{The publication chain strips metadata from evaluation
records. \textbf{Top:} a fully-specified record (model author, 9/9
fields) progressively loses configuration metadata as it passes through
the OLv2 leaderboard API, an aggregator, and a final comparison table.
The fields with zero valid coverage in our 10-source leaderboard subset
(\S\ref{sec:discrepancies}), \texttt{harness},
\texttt{scoring\_mode}, \texttt{prompt\_template}, \texttt{temperature}
 are dropped at the leaderboard step; \texttt{n\_shot} and
\texttt{chain\_of\_thought} are typically compressed away by the
comparison-table stage. \textbf{Bottom:} two stripped records meet on a
comparison table with exact 2-decimal agreement on Gemma-2-9B BBH
(34.10): one from OLv2, one from the SEA-LION paper
\citep{ng2025sealion}, which cites Open~LLM~Leaderboard among its
sources. Because the fields that would distinguish independent
replication from a re-imported score, \texttt{harness},
\texttt{prompt\_template\_hash}, \texttt{evaluation\_timestamp}
were dropped at step~2, the two records share a pipeline origin and cannot be treated as independent measurements (\S\ref{sec:sealion}).}
\label{fig:chain}
\end{figure*}

\subsection{The Misunderstanding Chain}

%% DRAFT --------------------------------------------------------------
% The findings of the audit and the case studies share a common shape.
% An evaluation result is produced by a long chain of decisions:
% which harness, which version, which prompt template, which n-shot
% count, which temperature, which scoring mode, which answer extractor.
% At evaluation time, every link in the chain is known: the runtime
% environment carries the harness version, the prompt template is a
% string in memory, the temperature is a float. Between evaluation
% time and the reader, almost all of this information is dropped.
% Leaderboard APIs return scores and model identifiers but rarely
% expose the configuration that produced them. Aggregators store
% whatever the APIs return. Downstream papers cite the aggregator.
% By the time two scores meet on a leaderboard page, the procedures
% behind them are no longer comparable in any verifiable way.
%
% Two alternatives to recording the configuration are sometimes
% suggested. The first is independent re-evaluation: if scores
% diverge, run the evaluation again under known conditions.
% Section~\ref{sec:results-discrepancies} shows that this is not
% available at scale. The model--benchmark overlap between sources is
% too sparse for cross-source verification to arise naturally, and
% targeted re-evaluation across thousands of models is not realistic.
% The second is automatic re-extraction from public artefacts. The
% Papers With Code artefact rate (57\% of raw records removed in our
% pipeline) shows that automatic ingestion of unstructured sources
% produces a substantial fraction of semantically wrong but
% schema-valid records, and the harness-identification audit
% (Section~\ref{sec:harness-audit}, pending) is expected to show that
% even relatively explicit information is hard to recover reliably.
%
% This leaves documentation. The argument is not that documentation is
% sufficient --- one can imagine adversarial cases where it is not ---
% but that in the current ecosystem it is necessary, because the two
% alternatives have failed. The closest analogy is the move to
% pre-registered clinical trials: when independent replication of every
% trial is infeasible, mandatory protocol recording becomes the
% substitute. Evaluation reporting is in a similar position, and
% requires a similar response.
%% END DRAFT ----------------------------------------------------------
An evaluation result is produced by a long chain of decisions: which
harness, which version, which prompt template, which n-shot count, which
temperature, which scoring mode, which answer extractor. At evaluation
time, every link in that chain is known. The runtime environment carries
the harness version, the prompt template is a string in memory and the
temperature is a float in a configuration file. Between evaluation time
and the reader, almost all of this information is dropped. Leaderboard
APIs return scores and model identifiers but rarely expose the
configuration that produced them. Aggregators store whatever the APIs
return. Downstream papers cite the aggregator. By the time two scores
meet on a comparison table, the procedures behind them are no longer
comparable in any verifiable way.

The three case studies in \S\ref{sec:cases} document where this loss
occurs: at the leaderboard (\S\ref{sec:protocol}, where OLv2 does not
expose harness or scoring mode), at the benchmark label
(\S\ref{sec:onenametwo}, where TruthfulQA, Belebele, and HellaSwag conceal
n-shot, metric variant, and scoring mode), and at the appearance of
independent measurement (\S\ref{sec:convergence}, where shared pipelines
produce apparent replication that is illusory).

Two alternatives to documentation are sometimes proposed. The first is
automatic re-extraction from public artefacts. Our extraction pipeline is
reliable when the source contains the information: the score-verification
audit (\S\ref{sec:score-ann}) confirms 100\% agreement across 1{,}619
row-level annotations on 10 papers, and the harness-identification audit
(\S\ref{sec:harness-ann}) achieves 86.2\% per-row precision on the
with-harness sample (3/5 papers at 100\%; the remaining 2 expose a
specific failure mode where broad group attributions like ``across all
aforementioned tasks'' are under-extracted) and 100\% correct abstention
on the without-harness sample. But Figure~\ref{fig:coverage} shows that
in the 10-source leaderboard subset, three of five configuration fields~---
harness identity, prompt template, and temperature~--- have zero valid
coverage.
No extraction effort can recover what the source never recorded. The
second alternative is independent re-evaluation, which is feasible at
small scale but not across the thousands of model-benchmark pairs in a
modern aggregated dataset. This leaves documentation: not sufficient in
every case, but necessary in the current public ecosystem, because the
two alternatives fail precisely where the need is greatest.

\subsection{EvalSpec v0.1 as a Proposed Remedy}
\label{sec:evalspec}

%% DRAFT --------------------------------------------------------------
% EvalSpec v0.1 is a machine-readable reporting specification for
% evaluation runs. It is designed to be the smallest set of fields that
% turns a score into an identifiable measurement: enough that two
% records meeting on the same benchmark can be compared without an
% intervening act of guesswork.
%
% Field selection follows from the audit. The Required tier
% (EvalSpec-R) contains the fields whose absence makes attribution
% impossible: model identifier, benchmark identifier, scoring mode
% (one of log-likelihood, generation, preference, execution), harness
% name and version, n-shot count, date, and random seed. The
% Recommended tier (EvalSpec-Rec) contains the fields whose absence
% makes attribution incomplete: prompt-template hash, temperature,
% top-p, quantisation, hardware, software environment. The Full tier
% (EvalSpec-F) contains the fields needed for full reproducibility:
% provenance link, contamination check, compute hours.
% Table~\ref{tab:evalspec-fields} lists the full set.
%
% EvalSpec is encoded as a JSON-LD vocabulary extending schema.org
% and compatible with Croissant metadata records. Compliance is
% enforced by SHACL shape constraints, so that a record can be
% mechanically checked against any tier. The prompt template is
% recorded as a SHA-256 hash of the canonical template string, which
% keeps the specification zero-overhead while preserving the ability
% to detect template changes across runs. A reference validator and
% the full vocabulary are released alongside the dataset.
%
% None of the Required fields require new instrumentation. Temperature
% is one float per run; the prompt template is one hash; the harness
% name and version are strings already present in the runtime
% environment. The barrier to adoption is not implementation cost but
% API exposure: the leaderboards already know these values, but their
% APIs do not return them.
%% END DRAFT ----------------------------------------------------------

EvalSpec v0.1 is a machine-readable reporting specification designed to be
the smallest field set that turns a score into an identifiable measurement:
enough that two records on the same benchmark can be compared without
guesswork about the procedure that produced them.

Field selection follows directly from the case studies in
\S\ref{sec:cases}. \texttt{scoring\_mode} (log-likelihood vs.\ free-form
generation) is motivated by the BBH label-collapse (\S\ref{sec:bbh}) and
the HellaSwag analysis (\S\ref{sec:hellaswag}). \texttt{harness\_name}
and \texttt{harness\_version} are required because the SEA-LION case
(\S\S\ref{sec:bbh}--\ref{sec:sealion}) shows that harness identity is
the fingerprint distinguishing genuinely independent measurement from a
shared-pipeline artefact. \texttt{n\_shot} is motivated by the TruthfulQA,
Belebele, and HellaSwag audits
(\S\S\ref{sec:truthfulqa}--\ref{sec:hellaswag}), where the n-shot count
shifts the evaluation regime but is absent from every aggregated source's
benchmark label. \texttt{benchmark\_variant} (recording sub-task keys such
as \texttt{mc1} vs.\ \texttt{mc2}) is motivated by TruthfulQA and Belebele
(\S\S\ref{sec:truthfulqa}--\ref{sec:belebele}), where the same name covers
metrics with opposite definitions of correctness. \texttt{judge\_model} is
required because TruthfulQA (\S\ref{sec:truthfulqa}) uses a fine-tuned
judge whose identity is material to the score and absent from the
aggregated record; a controlled experiment (Appendix~\ref{app:exp-judge})
confirms the effect is not marginal---varying only the judge over fixed
responses moves mean scores by up to 1.35 points on a 10-point scale and
reorders the system ranking (cross-judge rank correlations span $+1.0$ to
$-0.3$), so the judge is not an offset that cancels in a comparison. \texttt{evaluation\_timestamp} is required because the
Math-Verify parser replacement (\S\ref{sec:math}) demonstrates that a
silent infrastructure change produces a different score under the same
benchmark label. The full field set by tier is in
Table~\ref{tab:evalspec-fields-app} (Appendix~\ref{app:evalspec}).

EvalSpec is encoded as a JSON-LD vocabulary extending \texttt{schema.org}
and compatible with Croissant metadata records. Compliance is enforced by
SHACL shape constraints, so that any record can be mechanically checked
against a target tier without human judgement. The prompt template is
recorded as a SHA-256 hash of the canonical template string: this keeps the
specification zero-overhead for producers while preserving the ability to
detect template changes across runs or harness upgrades. A reference
validator and the full vocabulary are released alongside the dataset.

None of the Required fields require new instrumentation: the values exist
in the runtime environment at evaluation time. The barrier to adoption is
not implementation cost but API exposure~--- leaderboards already record
these values internally, but their public interfaces do not return them.


\subsection{Recommendations}

%% DRAFT --------------------------------------------------------------
% Concrete recommendations follow from the audit, the case studies,
% and the EvalSpec design.
%
% \paragraph{Leaderboard maintainers} should expose harness identity,
% harness version, prompt-template hash, and temperature in their
% APIs. These values are already known at evaluation time --- the
% leaderboard runs the evaluation --- and adding them to the API
% response requires only logging. The structural barrier here is
% incentive, not implementation.
%
% \paragraph{Model developers} should ship EvalSpec-R compliant
% evaluation records alongside model releases. A small set of
% well-documented records, even from a single evaluation pass, makes
% the model's reported scores verifiable against subsequent
% third-party evaluations.
%
% \paragraph{Schema and platform designers} should treat scoring mode
% as a first-class field. No combination of the other fields
% recovers it: a record can be fully populated on harness, prompt,
% temperature, and n-shot, and still leave the log-likelihood-versus-
% generation distinction implicit. The GPQA collision in
% Section~\ref{sec:case-gpqa} demonstrates the consequence of leaving
% this field implicit at the scale of the public ecosystem.
%% END DRAFT ----------------------------------------------------------

\paragraph{Leaderboard maintainers.}
Leaderboard operators should expose \texttt{harness\_name},
\texttt{harness\_version}, \texttt{scoring\_mode}, and
\texttt{evaluation\_timestamp} through their public APIs. These values are
already known at evaluation time; adding them is a logging decision, not
an engineering project. The SEA-LION case (\S\ref{sec:sealion}) shows the
cost of not exposing them: an exact score match becomes unreliable
evidence of independent replication. Records meeting the EvalSpec
Required tier would let downstream researchers separate shared-pipeline
artefacts from genuine replications.

\paragraph{Model developers.}
Model releases should include at least one EvalSpec-R compliant record per
reported benchmark. A single well-documented record makes scores verifiable
against subsequent third-party evaluations and provides a concrete baseline
for detecting benchmark variant drift. E1 (Appendix~\ref{app:exp-e1})
shows the magnitude of drift available from a single unrecorded prompt
format: 58.5\% vs.\ 2.5\% on the same model, benchmark, and harness. The
TruthfulQA and Belebele discrepancies
(\S\S\ref{sec:truthfulqa}--\ref{sec:belebele}) show the same conflation at
the level of n-shot, metric variant, and judge; a machine-readable record
anchors comparison to the configuration actually used in the release.

\paragraph{Schema and platform designers.}
\texttt{scoring\_mode} should be treated as a first-class
controlled-vocabulary field in any evaluation metadata schema. No
combination of other Required fields recovers it: a record can be fully
specified on harness, n-shot, and timestamp and still leave the
log-likelihood-vs.-generation boundary implicit. The BBH label-collapse
(\S\ref{sec:bbh}) shows the consequence at ecosystem scale, and E2
(Appendix~\ref{app:exp-e2}) shows the converse: when scoring mode is the
only varied factor, the gap collapses to under 0.3 pp. A four-value
controlled vocabulary (\texttt{log-likelihood}, \texttt{generation},
\texttt{preference}, \texttt{execution}) closes this gap.



\section{Conclusions}

%% DRAFT --------------------------------------------------------------
% Public evaluation records routinely report scores without the
% procedures that produced them. Across 29{,}331 records from eleven
% sources, the configuration fields most responsible for score
% divergence --- harness identity, prompt template, and temperature
% --- have zero valid coverage. The cross-source overlap needed for
% independent verification is rare, and concentrated in a single
% dominant source. Three case studies show how the absence of these
% fields makes specific divergences uninterpretable.
%
% Documentation is the only remaining identification strategy.
% EvalSpec v0.1 is a concrete proposal for what to record, grounded in
% the audit and validated against the case studies. Dataset, pipeline,
% schema, and analysis code are released under CC-BY-4.0.
%% END DRAFT ----------------------------------------------------------

\textit{[To be written last, once main sections settle.]}
\section*{Acknowledgments}
\bibliography{custom}

\appendix

\section{Score-Verification Audit Detail}
\label{app:audit}

Table~\ref{tab:score-ann-app} reports per-paper row counts for the
score-extraction accuracy audit (\S\ref{sec:score-ann}). All 1{,}619
entries verified at 100\% agreement.

\begin{table}[ht]
\centering
\small
\begin{tabular}{lrr}
\toprule
Paper & ArXiv ID & Entries \\
\midrule
Instella                 & 2511.10628 & 322 \\
Kimi-Audio               & 2504.18425 & 249 \\
Sailor                   & 2404.03608 & 205 \\
DeepSeek LLM             & 2401.02954 & 190 \\
Llama-Embed-Nemotron-8B  & 2511.07025 & 180 \\
OpenELM                  & 2404.14619 & 152 \\
Phi-4-reasoning-vision   & 2603.03975 & 150 \\
Mistral 7B               & 2310.06825 &  62 \\
LLM360                   & 2312.06550 &  56 \\
DataComp-LM              & 2406.11794 &  53 \\
\midrule
Total                    &            & 1{,}619 \\
\bottomrule
\end{tabular}
\caption{Per-paper row counts in the score-accuracy annotation set
(\S\ref{sec:score-ann}).}
\label{tab:score-ann-app}
\end{table}


\section{Harness-Identification Audit Methodology}
\label{app:harness}

For each of the 10 audit papers (\S\ref{sec:harness-ann}), the ArXiv HTML
was retrieved from \texttt{arxiv.org/html/} and the pipeline's per-row
\texttt{eval\_library.name} value was verified against the source text by
hand. The audit operates at the (model, benchmark) row level rather than
the paper level, because a single paper can use a named harness for some
benchmarks and a separate pipeline (machine-translation suite, code
sandbox, LLM-as-judge, proof assistant) for others.

\paragraph{Strict-match rubric.} A row is \textsc{true} only if the
pipeline's harness label correctly names the tool the paper says was used
for that specific (model, benchmark) cell. The rubric accepts canonical
equivalents for the same library: ``eval-harness'', ``lm-evaluation-harness'',
``Language Model Evaluation Harness'', and ``LM-Evaluation Harness'' all
refer to EleutherAI's library and are treated as interchangeable. A
paper's distinctly-named fork (e.g., BLOOM's \emph{Prompted Language
Model Evaluation Harness}, the promptsource-extended fork the BLOOM
authors named, built, and released) is required as a verbatim match: the
pipeline must produce the fork's name, not the parent library's name,
because the two tools can produce different scores on the same input.

\paragraph{Sample composition.} The \emph{with-harness} pool has five
papers (BLOOM, Llemma, Sheared LLaMA, Aya 23, Nemotron-4 15B) covering
817 rows; the \emph{without-harness} pool has five (Code Llama, Mamba,
Reka, Jamba, a Llama-2-related survey paper) covering 816 rows. The two
pools are disjoint from the score-verification sample
(Appendix~\ref{app:audit}). Per-row precision and abstention are reported
in \S\ref{sec:harness-ann}.


\section{Case-Study Score Tables}
\label{app:casetables}

\begin{table}[ht]
\centering
\small
\begin{tabular}{lll}
\toprule
Source & Score & Notes \\
\midrule
OLv2 (lighteval, log-likelihood) & 21.12 & 3-shot \\
Jamba paper & 55.10 & CoT generation \\
RecurrentGemma paper & 55.10 & CoT generation \\
MAmmoTH2 paper & 57.40 & CoT generation \\
(BigBench-Hard alias) & 59.60 & CoT generation \\
\bottomrule
\end{tabular}
\caption{Gemma-7B BBH scores across sources (\S\ref{sec:bbh}). The
$\approx$35-point gap reflects OLv2's lighteval log-likelihood scoring
vs.\ the ArXiv papers' chain-of-thought generation. }
\label{tab:bbh-app}
\end{table}

\begin{table}[ht]
\centering
\small
\begin{tabular}{lll}
\toprule
Source & Score & Notes \\
\midrule
OLv2 (lighteval, log-likelihood) & 34.10 & 3-shot \\
SEA-LION paper & 34.10 & shared pipeline origin \\
OLMoE paper & 68.20 & 0-shot CoT generation \\
SmolLM2 paper & 69.00 & 0-shot CoT generation \\
RecurrentGemma paper & 69.00 & 0-shot CoT generation \\
\bottomrule
\end{tabular}
\caption{Gemma-2-9B BBH scores across five sources (\S\ref{sec:sealion}).
Two clusters: shared pipeline origin (OLv2/SEA-LION) and protocol mismatch
(3-shot log-likelihood vs.\ 0-shot CoT).}
\label{tab:sealion-app}
\end{table}


\section{Metadata Coverage Table}
\label{app:coverage}

Table~\ref{tab:coverage-app} reports the numerical values underlying the
coverage bar chart in Figure~\ref{fig:coverage} (\S\ref{sec:discrepancies}).

\begin{table}[ht]
\centering
\small
\begin{tabular}{lcc}
\toprule
Field & Full dataset & 10-source subset \\
\midrule
shots & 86.7\% & 98.0\% \\
harness$^{\dagger}$ & 74.2\% & 93.2\% \\
chain\_of\_thought & 60.5\% & 77.7\% \\
prompt\_template & 7.0\% & 0.0\% \\
temperature & 1.1\% & 0.0\% \\
\bottomrule
\end{tabular}
\caption{Metadata coverage (fraction of records with a non-null,
non-``unknown'' value). Full dataset: 37{,}718 records / 12 sources.
10-source subset: 28{,}755 records, excluding the ArXiv extraction
source. Numerical values underlying Figure~\ref{fig:coverage}.
$^{\dagger}$The harness row counts pipeline-auto-attached values
(e.g., \texttt{lm\_eval} attached to every OLv2 record from platform
documentation, not returned by the source API). No leaderboard source
exposes harness identity directly.}
\label{tab:coverage-app}
\end{table}


\section{EvalSpec v0.1 Field Set}
\label{app:evalspec}

Table~\ref{tab:evalspec-fields-app} lists the complete EvalSpec v0.1 field
set by compliance tier, motivated in \S\ref{sec:evalspec}.

\begin{table}[!htbp]
\centering
\small
\setlength{\tabcolsep}{4pt}
\begin{tabular}{lll}
\toprule
\textbf{Tier} & \textbf{Field} & \textbf{Description} \\
\midrule
Required & \texttt{model\_id}               & Model identifier (HF path or equivalent) \\
         & \texttt{benchmark\_id}           & Benchmark name and version \\
         & \texttt{scoring\_mode}           & \{log-likelihood, generation, preference, execution\} \\
         & \texttt{harness\_name}           & Evaluation harness name \\
         & \texttt{harness\_version}        & Exact harness version string \\
         & \texttt{n\_shot}                 & Number of in-context examples \\
         & \texttt{benchmark\_variant}      & Sub-task or configuration key \\
         & \texttt{judge\_model}            & Scorer model (\texttt{null} if not applicable) \\
         & \texttt{evaluation\_timestamp}   & ISO-8601 date of evaluation run \\
\midrule
Recommended & \texttt{prompt\_template\_hash} & SHA-256 of canonical prompt template \\
            & \texttt{temperature}           & Sampling temperature \\
            & \texttt{top\_p}                & Top-$p$ truncation threshold \\
            & \texttt{quantisation}          & Weight quantisation scheme \\
            & \texttt{hardware}              & GPU type and count \\
            & \texttt{random\_seed}          & RNG seed \\
\midrule
Full & \texttt{provenance\_url}             & Link to raw evaluation log or artefact \\
     & \texttt{contamination\_check}        & Decontamination method and date \\
     & \texttt{compute\_hours}              & Wall-clock GPU-hours \\
\bottomrule
\end{tabular}
\caption{EvalSpec v0.1 field set by compliance tier (\S\ref{sec:evalspec}).
Required fields are the minimum for score comparability;
Recommended fields support reproducibility;
Full fields enable complete audit.}
\label{tab:evalspec-fields-app}
\end{table}


\FloatBarrier
\section{Additional Case-Study Evidence}
\label{app:hellaswag}

Figure~\ref{fig:unified-hellaswag} shows the unified case-study figure
with HellaSwag (\S\ref{sec:hellaswag}) as the middle column, an
alternative to the Belebele variant used in Figure~\ref{fig:casestudies-unified}.

\begin{figure*}[ht]
\centering
\includegraphics[width=\textwidth]{./figures/case_studies_unified_hellaswag_S3.pdf}
\caption{Alternative middle column for the unified figure
(cf.\ Figure~\ref{fig:casestudies-unified}), showing the HellaSwag case
study (\S\ref{sec:hellaswag}) instead of Belebele.
Phi-3 (5-shot, generation-based) vs.\ Gemma (0-shot, log-likelihood);
both prompts shown are canonical reconstructions since neither source
paper publishes its HellaSwag prompts verbatim.}
\label{fig:unified-hellaswag}
\end{figure*}
\FloatBarrier
\section{Experimental Setup}
\label{app:experiments}

Four experiments and one framing factorial grid anchor the analyses in
\S\ref{sec:cases}. Table~\ref{tab:experiments-master} summarises them;
per-experiment design and results follow in
\S\S\ref{app:exp-grid}--\ref{app:exp-judge}.

\begin{table*}[ht]
\centering
\small
\setlength{\tabcolsep}{4pt}
\begin{tabular}{lp{3.6cm}p{3.0cm}p{4.4cm}p{4.0cm}}
\toprule
\textbf{ID} & \textbf{Goal} & \textbf{Models} & \textbf{Setup} & \textbf{Result} \\
\midrule
Factorial Grid &
Quantify relative importance of methodological factors &
Qwen2.5-14B, Qwen2.5-7B, Llama-3.1-8B, Mistral-7B-v0.3 &
3 benchmarks $\times$ 3 formats $\times$ 2 temps $\times$ 2 n-shots $\times$ 3 seeds (MMLU reduced to 1 temp $\times$ 1 seed); 312 attempted, 306 completed &
Prompt format $\approx$27\% of variance on GSM8K ($F{=}13.0$, $p{<}0.001$); see \S\ref{app:exp-grid} for the full per-benchmark breakdown \\
\midrule
E1 &
Test parser-failure mechanism under varied prompt format &
Qwen2.5-14B-Instruct &
GSM8K, 5-shot, lm-evaluation-harness 0.4.11; prompt formats plain / instruct / CoT; temperature 0.0; seeds 42, 123, 7; 200-problem subset &
Plain 58.5\%, Instruct 57.0\%, CoT 2.5\% (\S\ref{sec:gsm8k}) \\
\midrule
E2 &
Isolate scoring-mode effect from other factors &
Mistral-7B-Instruct-v0.3, Qwen2.5-7B-Instruct, Llama-3.1-8B-Instruct (paired); Qwen2.5-14B-Instruct (generation only) &
MMLU, 5-shot, plain prompt, temperature 0.0; log-likelihood vs.\ generation-based scoring &
Maximum gap across three paired models: 0.28 pp (\S\ref{sec:bbh}, ruling out a generic scoring-mode effect) \\
\midrule
E3 &
Test whether the harness \emph{version} alone shifts scores &
Qwen2.5-1.5B-Instruct, Qwen2.5-7B-Instruct &
GSM8K full, 5-shot, \texttt{hf} backend; lm-evaluation-harness 0.4.3 vs.\ 0.4.11; seeds 42, 123 &
Cross-version $\le$0.45 pp, below the seed spread ($\le$1.29 pp); backs \texttt{harness\_version} + the SEA-LION argument (\S\ref{sec:sealion}) \\
\midrule
Judge &
Isolate the judge-model effect (LLM-as-judge) &
5 contestant models, 4 judges (GPT-4o-mini, Gemini-3.5-Flash, Llama-3.1-70B, Claude-Haiku) &
MT-Bench turn-1; identical responses scored by each judge &
Judge-induced spread up to 1.35 pts/10; ranking flips (cross-judge Spearman $+1.0$ to $-0.3$); see \S\ref{app:exp-judge} \\
\bottomrule
\end{tabular}
\caption{All experiments anchoring the analyses in \S\ref{sec:cases}. The
Factorial Grid is the framing experiment; E1, E2, E3, and the judge-model
experiment each isolate a specific mechanism. Per-cell results and full
configurations are in the subsections below.}
\label{tab:experiments-master}
\end{table*}


\subsection{Factorial Grid}
\label{app:exp-grid}

The framing experiment for \S\ref{sec:cases}. Four instruction-tuned
models (Qwen2.5-14B-Instruct, Qwen2.5-7B-Instruct, Llama-3.1-8B-Instruct,
Mistral-7B-Instruct-v0.3) were evaluated on three benchmarks (BBH, GSM8K,
MMLU) under three prompt formats (\texttt{plain} few-shot,
\texttt{instruct} system-prompt-formatted, \texttt{cot} chain-of-thought
few-shot), two temperatures (0.0, 0.7), and two n-shot counts per
benchmark: 0-shot plus 3-shot for BBH, 0-shot plus 5-shot for GSM8K and
MMLU. Three random seeds (42, 123, 7) were used per cell. MMLU was
reduced to temperature 0.0 with a single seed (42) due to
evaluation-server constraints. The analysis covers 306 successful evaluations:
144 BBH cells (4 models $\times$ 2 n-shots $\times$ 3 formats $\times$ 2
temperatures $\times$ 3 seeds), 144 GSM8K cells (same design), and 18
MMLU cells (3 models $\times$ 2 n-shots $\times$ 3 formats at temperature
0.0 with a single seed).

Inference was run with lm-evaluation-harness 0.4.11 and vLLM. GSM8K and
MMLU runs used a 200-problem subset (\texttt{limit=200}) for compute
budget; BBH used the full task suite. Variance attribution uses two
separate one-way ANOVAs on the prompt-format factor, computed per
benchmark on the 5-shot cells to avoid pooling across the unbalanced
GSM8K and MMLU designs. GSM8K contributes 72 cells (4 models $\times$ 3
formats $\times$ 2 temperatures $\times$ 3 seeds): prompt format yields
$F = 13.00$, $p = 1.6 \times 10^{-5}$, $\eta^2 \approx 27.4\%$, while
temperature contributes $\eta^2 \approx 0.25\%$ ($F = 0.18$, $p = 0.68$).
MMLU's reduced 9-cell design at 5-shot (3 models $\times$ 3 formats
$\times$ 1 temperature $\times$ 1 seed) is descriptively consistent
($\eta^2 \approx 24.5\%$) but underpowered ($F = 0.98$, $p = 0.43$); the
per-format means agree with the GSM8K pattern but cannot independently
confirm it. Per-format means for both benchmarks are reported in
Table~\ref{tab:anova-per-benchmark}.

We attribute variance to the n-shot factor the same way, pooling each
benchmark's successful cells by shot count (0-shot vs.\ 5-shot;
Table~\ref{tab:anova-nshot}), and the two benchmarks behave oppositely in a
way that is itself diagnostic. On MMLU the n-shot effect is negligible:
pooled across models and formats, 0-shot scores 61.94 and 5-shot 59.34
(\(\eta^2 \approx 1.0\%\), \(F = 0.16\), \(p = 0.69\)). This corroborates
the HellaSwag analysis (\S\ref{sec:hellaswag}) with a controlled
measurement: on a log-likelihood-scored benchmark, varying n-shot moves the
score by at most a few points. On GSM8K the n-shot factor appears dominant
(\(\eta^2 \approx 65\%\) pooled, \(80\%\) within \texttt{plain}), but the
effect is not a capability gain: every 0-shot cell scores exactly 0.0,
because the instruction-tuned models without few-shot priming do not emit
the \texttt{\#\#\#\#} answer delimiter the GSM8K extractor requires, so the
parser records zero for every item. This is the same parser--format
coupling documented in Case Study~1 (\S\ref{sec:gsm8k}), here triggered by
shot count rather than prompt format: on a generation-and-parse benchmark
n-shot and scoring interact, whereas on a log-likelihood benchmark n-shot is
nearly inert.

\begin{table}[h]
\centering
\small
\begin{tabular}{llrrr}
\toprule
Benchmark & Format & $n$ & Mean & Range \\
\midrule
GSM8K & \texttt{plain}    & 24 & 38.75 & 18.50--59.50 \\
GSM8K & \texttt{instruct} & 24 & 34.19 & 16.50--57.00 \\
GSM8K & \texttt{cot}      & 24 & 19.25 & ~~2.00--41.00 \\
\midrule
MMLU  & \texttt{plain}    &  3 & 63.60 & 55.17--73.18 \\
MMLU  & \texttt{instruct} &  3 & 62.45 & 52.52--72.46 \\
MMLU  & \texttt{cot}      &  3 & 51.98 & 35.97--61.95 \\
\bottomrule
\end{tabular}
\caption{Per-benchmark, per-format score distributions at 5-shot in the
factorial grid. GSM8K cells are pooled across 4 models $\times$ 2
temperatures $\times$ 3 seeds; MMLU cells across 3 models $\times$ 1
temperature $\times$ 1 seed. The plain/instruct/cot ordering of means is
the same on both benchmarks; the GSM8K plain--cot gap is 19.50 pp and
the MMLU plain--cot gap is 11.61 pp. Raw values released with the
analysis code.}
\label{tab:anova-per-benchmark}
\end{table}

\begin{table}[h]
\centering
\small
\begin{tabular}{llrr}
\toprule
Benchmark & n-shot & $n$ & Mean \\
\midrule
GSM8K & 0-shot & 72 & ~~0.00 \\
GSM8K & 5-shot & 72 & 30.73 \\
\midrule
MMLU  & 0-shot &  9 & 61.94 \\
MMLU  & 5-shot &  9 & 59.34 \\
\bottomrule
\end{tabular}
\caption{Per-benchmark, per-shot score means, pooled across models, prompt
formats, temperatures, and seeds. One-way ANOVA on the n-shot factor:
GSM8K \(\eta^2 \approx 65.1\%\) (\(F = 264.6\), \(p < 10^{-30}\)), MMLU
\(\eta^2 \approx 1.0\%\) (\(F = 0.16\), \(p = 0.69\)). The GSM8K 0-shot
mean of exactly 0.00 reflects parser failure, not absent capability: the
instruction-tuned models omit the \texttt{\#\#\#\#} answer delimiter
without few-shot priming (\S\ref{sec:gsm8k}), so the apparent n-shot effect
is the same parser--format coupling as Case Study~1. On MMLU
(log-likelihood) the n-shot effect is negligible.}
\label{tab:anova-nshot}
\end{table}



\subsection{E1: GSM8K Prompt-Format Sensitivity}
\label{app:exp-e1}

A slice of the Factorial Grid isolating the Qwen2.5-14B-Instruct $\times$
GSM8K $\times$ 5-shot cells at temperature 0.0. Three prompt formats are
compared with seeds 42, 123, 7 (nine runs total):

\begin{table}[ht]
\centering
\small
\begin{tabular}{lr}
\toprule
Prompt format (temp 0.0, seeds 42/123/7) & Score \\
\midrule
\texttt{plain} (direct-answer few-shot)      & 58.5\% \\
\texttt{instruct} (system-prompt-formatted)  & 57.0\% \\
\texttt{cot} (chain-of-thought few-shot)     & ~~2.5\% \\
\bottomrule
\end{tabular}
\caption{Qwen2.5-14B-Instruct on GSM8K, 5-shot, lm-evaluation-harness
0.4.11, \texttt{limit=200}. Scores at temperature 0.0 are identical
across all three seeds (deterministic decoding), so each row reports a
single number rather than a mean with variance. The \texttt{cot}
collapse to 2.5\% reflects parser failure: the standard GSM8K extractor
\texttt{re.search(r'\#\#\#\#\textbackslash s+(\textbackslash S+)', output)}
finds no delimiter in the chain-of-thought free-text output
(\S\ref{sec:gsm8k}).}
\label{tab:e1-gsm8k}
\end{table}

The same nine cells run at temperature 0.7 give means across three seeds
of 58.67\% (plain), 52.67\% (instruct), and 2.83\% (CoT), confirming that
the parser-failure mechanism holds across sampling regimes.


\subsection{E2: MMLU Scoring-Mode Control}
\label{app:exp-e2}

Three instruction-tuned models were evaluated on MMLU 5-shot under both
log-likelihood and generation-based scoring to test whether scoring mode
alone produces large score gaps. A fourth model (Qwen2.5-14B-Instruct)
contributes a generation-only run for reference.

\begin{table}[ht]
\centering
\small
\setlength{\tabcolsep}{4pt}
\begin{tabular}{lrrr}
\toprule
Model & log-likelihood & generation & $|\Delta|$ (pp) \\
\midrule
Mistral-7B-Instruct-v0.3   & 61.92 & 61.64 & 0.28 \\
Qwen2.5-7B-Instruct        & 74.26 & 74.43 & 0.17 \\
Llama-3.1-8B-Instruct      & 68.30 & 68.16 & 0.14 \\
Qwen2.5-14B-Instruct       &  ---  & 79.79 & ---  \\
\bottomrule
\end{tabular}
\caption{MMLU 5-shot, plain prompt, temperature 0.0, lm-evaluation-harness
0.4.11. Log-likelihood runs use a 500-question subset
(\texttt{limit=500}); generation runs use the full MMLU test set
(\texttt{limit=14042}). The maximum gap across the three paired models
is 0.28 pp, which rules out a generic scoring-mode effect: the
$\approx$35-point BBH gap documented in \S\ref{sec:bbh} is specific to
the interaction between lighteval's log-likelihood approach and BBH's
chain-of-thought structure, not a property of log-likelihood scoring in
general.}
\label{tab:e2-scoring-mode}
\end{table}


\subsection{E3: Harness-Version Sensitivity}
\label{app:exp-e3}

To isolate the effect of the evaluation-harness \emph{version}, two
instruction-tuned models (Qwen2.5-1.5B-Instruct, Qwen2.5-7B-Instruct) were
evaluated on the full GSM8K test set (5-shot, \texttt{hf} backend) under two
versions of lm-evaluation-harness, 0.4.3 and 0.4.11, with two random seeds
(42, 123). The model, benchmark, and all generation settings are held fixed, so
the only variable is the harness version (Table~\ref{tab:e3-harness-version}).

\begin{table}[ht]
\centering
\small
\begin{tabular}{llrr}
\toprule
Model & Harness & seed 42 & seed 123 \\
\midrule
Qwen2.5-1.5B-Instruct & 0.4.3  & 32.83 & 33.74 \\
Qwen2.5-1.5B-Instruct & 0.4.11 & 33.06 & 33.36 \\
Qwen2.5-7B-Instruct   & 0.4.3  & 75.06 & 76.35 \\
Qwen2.5-7B-Instruct   & 0.4.11 & 75.51 & ---   \\
\bottomrule
\end{tabular}
\caption{GSM8K (full test set, 5-shot, \texttt{hf} backend) under
lm-evaluation-harness 0.4.3 vs.\ 0.4.11 with identical configuration. The
largest seed-matched cross-version difference is 0.45 pp (Qwen2.5-7B, seed 42),
below the within-version seed spread (up to 1.29 pp on Qwen2.5-7B/0.4.3). The
0.4.11 / 7B seed-123 cell is omitted. Raw records released with the analysis
code.}
\label{tab:e3-harness-version}
\end{table}

The harness version shifts the GSM8K score by less than the run-to-run seed
variation: cross-version differences are at most 0.45 pp, while seeds alone move
the same configuration by up to 1.29 pp. The effect is real but small relative
to seed noise, which is precisely why \texttt{harness\_version} belongs in the
EvalSpec Required tier (\S\ref{sec:evalspec})---so that a future divergence can
be attributed to the version rather than guessed at. The same control bounds the
run-to-run variability invoked in the SEA-LION analysis (\S\ref{sec:sealion}).


\subsection{Judge-Model Sensitivity}
\label{app:exp-judge}

Motivated by the TruthfulQA judge confound (\S\ref{sec:truthfulqa}) and the four
LLM-judge sources in the dataset (Chatbot Arena, AlpacaEval~2, MT-Bench,
WildBench), this experiment isolates the effect of the judge model. Five
instruction-tuned ``contestant'' models each generated one response to the 80
turn-1 MT-Bench questions \citep{zheng2023mtbench}; the identical responses were
then scored by four judge models from different providers (GPT-4o-mini,
Gemini-3.5-Flash, Llama-3.1-70B-Instruct, Claude-Haiku-4.5) using the canonical
MT-Bench single-answer 1--10 grading rubric at temperature~0. Because the
responses are held fixed, any variation in a contestant's mean score across
judges is attributable to the judge alone. Each cell aggregates 78--80 ratings,
except Gemini-3.5-Flash (48--63), which intermittently returned empty completions.

Table~\ref{tab:judge-sensitivity} reports the mean rating per
(contestant, judge). The judge-induced spread---the gap between a contestant's
highest and lowest judge mean---averages 1.08 points and reaches 1.35 points on
the 10-point scale. More consequentially, the judge choice reorders the
leaderboard: Ministral-8B is ranked first by GPT-4o-mini and Llama-3.1-70B but
last by Gemini-3.5-Flash. Pairwise Spearman rank correlations between judges
span the full usable range, from $+1.00$ (GPT-4o-mini vs.\ Llama-3.1-70B,
identical order) down to $-0.30$ (GPT-4o-mini vs.\ Gemini-3.5-Flash,
anti-correlated). The judge model is therefore not a calibration offset that
cancels in a comparison; it can invert which system appears better, so a
benchmark score produced by an LLM judge is uninterpretable without recording
the judge (\S\ref{sec:evalspec}).

\begin{table}[ht]
\centering
\small
\setlength{\tabcolsep}{4pt}
\begin{tabular}{lrrrrr}
\toprule
Contestant & GPT-4o-m & Gem-3.5 & Llama-70B & Cl-Haiku & spread \\
\midrule
Gemma-3-4B   & 8.24 & 7.96 & 8.31 & 7.26 & 1.04 \\
Qwen2.5-7B   & 7.86 & 8.02 & 7.96 & 6.97 & 1.05 \\
Llama-3.1-8B & 8.04 & 7.61 & 8.22 & 6.87 & 1.35 \\
Ministral-8B & 8.44 & 7.54 & 8.62 & 7.51 & 1.11 \\
Qwen2.5-72B  & 8.32 & 8.59 & 8.38 & 7.74 & 0.84 \\
\bottomrule
\end{tabular}
\caption{Mean MT-Bench turn-1 rating (1--10) per contestant under each judge;
identical responses are scored by all four judges. ``spread'' is the
max$-$min judge mean for that contestant (mean 1.08, max 1.35). Ministral-8B
ranks first under GPT-4o-mini and Llama-70B but last under Gemini; cross-judge
Spearman rank correlation ranges from $+1.00$ to $-0.30$. Raw ratings released
with the analysis code.}
\label{tab:judge-sensitivity}
\end{table}


\end{document}
